\section{Introduction}
\label{sec:introduction}
% ============================================================
% MARKING CRITERION: "Abstract, Aims, Objectives and Project Description" [10 marks]
%   - Research aims of project are clear

%\guidance{%
%  MARKING: The introduction also contributes to the ``Aims, Objectives and
%  Project Description'' criterion (10 marks). Make the research aims and
%  objectives of the project unmistakably clear. The marker will look for:
%  (a) a well-motivated problem statement, (b) explicit research aims /
%  objectives, (c) a summary of the approach taken, and (d) a roadmap of
%  the report.}

\subsection{Problem Statement and Motivation}
While LLM-based educational agents offer high linguistic fluency, 
most existing systems rely on dyadic (one-on-one) text interactions that fail to capture the social dynamics of classroom learning. 
A key pedagogical gap lies in the absence of vicarious learning, the process where students gain knowledge by observing interactions between others. 
Passive learners often struggle to identify their own knowledge gaps, a problem exacerbated by the cognitive load of managing a dialogue. 
By moving away from text-only interfaces toward a triadic, real-time spoken "fictive dialogue", 
we can simulate a more natural and effective learning environment that triggers deeper reflection through agentic exchanges.

% Describe the broader context and why this problem matters.
% What gap or need does your work address?

\subsection{Research Aims and Objectives}
The primary aim of this project is to demonstrate how a multi-agent architecture can foster vicarious learning through autonomous student-teacher interactions.
The specific objectives are:
\begin{enumerate}[noitemsep, leftmargin=*]
  \item To develop a Teacher Agent capable of delivering structured, pedagogically sound content via a spoken channel.
  \item To implement an autonomous Student Agent that mimics human curiosity by interjecting with clarifying questions.
  \item To integrate low-latency speech technologies (Whisper STT and Chatterbox TTS) to manage a fluid triadic interaction.
  \item To evaluate the impact of a "synthetic peer" on user engagement and conceptual retention through controlled human testing.
\end{enumerate}

%\guidance{%
%  State your aims and objectives explicitly. Consider using a short
%  enumerated list so the marker can see them at a glance.}

% Example:
% The aims of this project are:
% \begin{enumerate}[noitemsep]
%   \item To implement ...
%   \item To evaluate ...
%   \item To investigate ...
% \end{enumerate}

\subsection{Experimental overview}

To demonstrate the effectiveness of our hypothesis, we conducted a user study comparing two pedagogical conditions: 
a standard dyadic interaction (Teacher-Human) and our proposed triadic fictive dialog (Teacher-Student-Human).

Participants are tasked with learning a specific topic through a spoken session in real-time. 
We measure learning gains through pre- and post-interaction assessments, while qualitative engagement is captured through subjective questionnaires. 
This setup allows us to isolate the impact of the Student Agent’s interjections on the user’s understanding and overall experience.
